\documentclass[11pt,a4paper]{article}
\usepackage[utf8]{inputenc}
\usepackage{amsmath,amssymb}
\usepackage{booktabs}
\usepackage[margin=2.5cm]{geometry}
\usepackage{hyperref}
\hypersetup{hidelinks}

\title{The $a_0$--$H_0$ relation in $\Lambda$-EG:\\
a consistency test and its distance-ladder circularity}

\author{J.\,P.\,Figueroa Torres\\
\small Independent Researcher --- ORCID: 0009-0005-5297-8777}

\date{August 2026 --- v2}

\begin{document}
\maketitle

\begin{abstract}
The Lambda Emergent Gravity ($\Lambda$-EG) framework fixes two acceleration
scales in a rigid geometric relation: the screening scale $g_{\rm crit} =
cH_0/(4\pi^2\sqrt{3})$ derived from de Sitter spectral geometry [Figueroa
Torres 2026], and the MOND acceleration $a_0 = 4\pi\, g_{\rm crit}$ derived from
the AQUAL embedding [Figueroa Torres 2026b, \S4]. We write out the immediate
cosmological consequence of these two identities, previously distributed
across companion papers but never stated explicitly: $H_0 = \pi\sqrt{3}\,
a_0/c$. Substituting the SPARC-canonical $a_0 = (1.20 \pm 0.02_{\rm stat}
\pm 0.24_{\rm sys})\times 10^{-10}\,\rm m\,s^{-2}$ yields $H_0 = 67.2$
km\,s$^{-1}$\,Mpc$^{-1}$, within $0.3\%$ of the Planck 2018 CMB central
value. We then show --- and this is the main cautionary result of the letter
--- that this agreement cannot at present be read as an independent
determination: between $55\%$ and $74\%$ of the standard SPARC sample
(depending on how cluster-assigned distances are counted) carries distances
computed under an assumed $H_0 = 73$ km\,s$^{-1}$\,Mpc$^{-1}$, and because
$a_0 \propto D^{-2}$ these propagate into $a_0$ with sample-weighted
sensitivity $\mathrm{d}\ln a_0/\mathrm{d}\ln H_0 = 1.1$--$1.5$, exceeding
unity and comparable to the tension being probed. Imposing the relation and
the calibration simultaneously yields a loop that is singular at
$H_0 \simeq 72$ and therefore determines no value of $H_0$ at all. We
present the relation as a consistency condition, quantify the coupling
explicitly, and identify the clean route: the $33$ SPARC galaxies with
individually measured redshift-independent distances.
We then observe that the companion
observable $\sigma_\Phi(z)$, the intrinsic scatter of the baryonic Tully--Fisher
relation across redshift, measures the shape $E(z) \equiv H(z)/H_0$
independently of $H_0$ itself, since $\sigma_\Phi(z)/\sigma_\Phi(0) =
E(z)^{\beta'(0)}$ [Figueroa Torres 2026b, \S4 and Appendix A.11]. The two
observables are informationally orthogonal: $a_0$ carries the scale, $\sigma_\Phi$
carries the shape. Combined, they reconstruct $H(z) = H_0\, E(z)$ from galactic
dynamics alone, using neither the CMB nor a distance ladder. We register
predictions before ELT/JWST-IFU deployment. All algebraic results follow from
previously published derivations; the new content is their combination and
its empirical evaluation.
\end{abstract}

\medskip\noindent\textbf{Keywords:} Hubble tension; MOND; emergent gravity;
baryonic Tully--Fisher relation; distance-ladder systematics

\section{Introduction}

The Hubble tension --- the persistent, now $\gtrsim 5\sigma$ discrepancy
between the value of $H_0$ inferred from the cosmic microwave background
under $\Lambda$CDM [Planck Collaboration 2020] and the value measured
directly through the Cepheid--supernova distance ladder [Riess et al. 2024]
--- has grown rather than shrunk with each generation of data. JWST
observations of the Cepheid rungs have ruled out crowding and dust as the
primary systematic in the local rung [Riess et al. 2024], and the two
determinations remain approximately $6$ km\,s$^{-1}$\,Mpc$^{-1}$ apart. The
tension therefore now motivates either an unrecognized systematic in one of
the two anchors, or a modification of the cosmological model between the
recombination epoch and today.

A third measurement path, independent of both anchors, would clarify which
possibility to take seriously. This letter identifies such a path within
$\Lambda$-EG, using observables that live entirely on galactic scales and
therefore share no calibration systematic with either the CMB or the distance
ladder. The path splits naturally into two orthogonal components: a
present-epoch determination of $H_0$ from the MOND acceleration scale $a_0$,
and a redshift-evolution determination of the shape $E(z) = H(z)/H_0$ from
the intrinsic scatter of the baryonic Tully--Fisher relation (BTFR). Neither
is new mathematics: the ingredients have all been derived and published in
the two companion papers [Figueroa Torres 2026, 2026b]. What is new is their
explicit combination, and its empirical evaluation.

We are honest at the outset about the strength of the present claim. The
determination of $H_0$ from $a_0$ is currently limited by the $\sim 20\%$
systematic on $a_0$ from stellar-mass modelling and inclination corrections
in SPARC [Lelli, McGaugh \& Schombert 2016]; the resulting band on $H_0$ is
wide enough to encompass both Planck and SH0ES. It is further limited, as we
show in \S4, by the fact that a majority of the SPARC distances are
themselves computed under an assumed $H_0$, so that the present-epoch channel
is not yet ladder-independent in the strict sense. What $\Lambda$-EG delivers
today is therefore neither a statistical arbitration of the tension nor a
fully independent third measurement, but a structural relation between two
acceleration scales that is testable, a quantified account of the
circularity that currently limits it, and a program to remove that
circularity through improved $a_0$ constraints and through ELT/JWST-IFU
measurements of $\sigma_\Phi(z)$. The second channel, $\sigma_\Phi(z)$, is
immune to the problem: being a ratio, it is independent of $H_0$ by
construction.

\section{The $a_0$--$H_0$ relation}

$\Lambda$-EG postulates two acceleration scales that are related by a fixed
geometric factor. The screening scale $g_{\rm crit}$ was derived in the
framework paper from de Sitter spectral geometry as
\begin{equation}
g_{\rm crit} \;=\; \frac{c H_0}{4\pi^2 \sqrt{3}}
\;\approx\; 9.58 \times 10^{-12}\ {\rm m\,s^{-2}},
\qquad H_0 = 67.4\,{\rm km\,s^{-1}\,Mpc^{-1}},
\label{eq:gcrit}
\end{equation}
[Figueroa Torres 2026, \S3]. The AQUAL embedding [Bekenstein \& Milgrom 1984]
of the same framework then fixes the deep-MOND acceleration scale $a_0$ that
enters the interpolating-function argument $x = g/a_0$ as
\begin{equation}
a_0 \;=\; 4\pi\, g_{\rm crit},
\label{eq:o6}
\end{equation}
where the factor $4\pi$ is the solid-angle weight of the Gauss-theorem flux
integration over a closed surface enclosing a spherically symmetric baryonic
source in the deep-MOND regime [Figueroa Torres 2026b, \S4]. Substituting
(\ref{eq:gcrit}) into (\ref{eq:o6}) and solving for $H_0$:
\begin{equation}
\boxed{\; H_0 \;=\; \frac{\pi \sqrt{3}\, a_0}{c} \;}
\label{eq:law}
\end{equation}
The prefactor $\pi\sqrt{3} \simeq 5.4414$ is a pure geometric constant. The
$\sqrt{3}$ is the square root of the spectral gap of the Laplacian on the
$S^3$ sections of de Sitter, $\lambda_1(S^3) = 3$, which the framework paper
takes as the primary origin of the factor [Figueroa Torres 2026, \S4.4]. The
identity $\lambda_1(S^D) = D$ makes the count independent of $D$ across
round spheres, and the framework argues separately that the same $\sqrt{D}$
follows from counting $D$ isotropic decoherence channels, which does not
presuppose a spherical section. We avoid the stronger phrasing
``topology-independent'': $\lambda_1 = D$ is a property of constant-positive-curvature
space forms, not of arbitrary topologies --- on a flat $T^3$ the gap depends
on the torus moduli, on $H^3$ it is set by the curvature radius, and on
$\mathbb{R}^3$ the spectrum is continuous with no gap at all. The claim that
survives is the weaker one: within space forms of constant positive
curvature the factor is fixed by the spatial dimension alone. The $\pi$
comes from the closed-surface integration of (\ref{eq:o6}).

Equation (\ref{eq:law}) is not a new derivation. It is the direct
consequence, in a single line of algebra, of (\ref{eq:gcrit}) and
(\ref{eq:o6}), both previously published. Its cosmological content lies in
the recognition that $a_0$ is measured on galactic scales, through
observables (rotation velocities, photometric masses) that involve no
cosmological model, so that (\ref{eq:law}) delivers $H_0$ without invoking
the CMB or the distance ladder \emph{as theoretical inputs}. Whether the
distances used to convert those observables into physical units are
themselves free of $H_0$ is a separate question, and \S4 shows that at
present they are not.

Two remarks. First, the factor $4\pi$ in (\ref{eq:o6}) is not a phenomenological
fit to SPARC: it is the Gauss-theorem projection weight, structural to any
scalar-field theory whose deep-MOND behaviour reduces to AQUAL. The observed
agreement with SPARC at the $\sim 4\%$ level is therefore data-limited rather
than fit-limited. Second, (\ref{eq:law}) is a hard prediction of the
combined $\Lambda$-EG + AQUAL structure: it cannot be relaxed by choosing a
different interpolating function, since $g_{\rm crit}$ and $a_0$ are both
defined outside the interpolant, at the asymptotic Newtonian and deep-MOND
boundaries respectively.

\section{Empirical evaluation}

The SPARC survey [Lelli, McGaugh \& Schombert 2016; McGaugh, Lelli \&
Schombert 2016] measures $a_0$ from the radial-acceleration relation across
approximately 150 disk galaxies, obtaining the canonical value
\begin{equation}
a_0 \;=\; (1.20 \pm 0.02_{\rm stat} \pm 0.24_{\rm sys})
\times 10^{-10}\ {\rm m\,s^{-2}},
\label{eq:a0sparc}
\end{equation}
where the systematic uncertainty is dominated by stellar mass-to-light
modelling and inclination corrections. Substituting the central value into
(\ref{eq:law}),
\begin{equation}
H_0 \;=\; \frac{\pi\sqrt{3}\times 1.20\times 10^{-10}}{2.998\times 10^{8}}\
{\rm s^{-1}}
\;=\; 67.21\ {\rm km\,s^{-1}\,Mpc^{-1}}.
\label{eq:H0central}
\end{equation}
Propagating the SPARC systematic on $a_0$ directly through (\ref{eq:law})
gives the interval $H_0 \in [53.8, 80.6]$ km\,s$^{-1}$\,Mpc$^{-1}$. Both
reference determinations sit comfortably inside it: Planck ($67.4 \pm 0.5$)
at $0.01\sigma_{\rm sys}$ and SH0ES ($73.5 \pm 0.7$) at $0.47\sigma_{\rm sys}$
from the central value. We stress that the band does not discriminate between
them at all --- the proximity of the central value to Planck carries no
statistical weight on its own. The comparison is summarised in
Table~\ref{tab:H0}.

\begin{table}[h]
\centering
\begin{tabular}{lcl}
\toprule
Method & $H_0$ [km\,s$^{-1}$\,Mpc$^{-1}$] & Anchor \\
\midrule
Planck 2018 CMB          & $67.4 \pm 0.5$   & recombination-epoch acoustic scale \\
SH0ES 2024               & $73.5 \pm 0.7$   & Cepheid--SN Ia distance ladder \\
\textbf{This letter} (Eq.~\ref{eq:law}) & $\mathbf{67.2}$ (central) & \textbf{galactic $a_0$ via SPARC} \\
\midrule
$\Delta$ vs.\ Planck  & $-0.3\%$ (central) & \\
$\Delta$ vs.\ SH0ES   & $-8.6\%$ (central) & \\
\bottomrule
\end{tabular}
\caption{Comparison of $H_0$ determinations from three independent anchors.
The $\Lambda$-EG prediction (this letter) uses no CMB data and no distance
ladder; its input is the SPARC-canonical value of $a_0$. The wide statistical
band from the SPARC systematic ($\pm 20\%$ on $a_0$ translates to $\pm 13.4$
km\,s$^{-1}$\,Mpc$^{-1}$ on $H_0$) prevents statistical arbitration at this
data quality; the coincidence of central values with Planck is nevertheless
notable given the total independence of the two anchors.}
\label{tab:H0}
\end{table}

The central value coincides with the Planck determination at a fractional
precision of $3 \times 10^{-3}$; the same central value differs from SH0ES
by $8.6\%$. Given that a shift of the $\Lambda$-EG central value from
$67.21$ to the SH0ES central of $73.5$ would require $a_0$ to have a true
value of $\sim 1.31 \times 10^{-10}$ m\,s$^{-2}$ (a $9\%$ excess over the SPARC
central), and given the empirical precision expected from Gaia-anchored
extended SPARC and LSST-era radial-acceleration-relation determinations
[order-of-magnitude reduction in the $a_0$ systematic on a decade timescale],
Eq.~(\ref{eq:law}) has the potential to become a sharp test. That potential
is conditional on removing the circularity quantified in \S4: a tighter
$a_0$ derived from the same distance mix would be a tighter number, not a
more independent one.

\section{The $H_0$ content of the $a_0$ measurement}
\label{sec:selfconsistency}

The argument of \S3 is only as independent of the distance ladder as the
distances that enter $a_0$. We examined this directly and found that it is
not independent at present data quality. The result is reported here rather
than in a footnote because it bounds the claim the letter can make.

Table~1 of SPARC [Lelli, McGaugh \& Schombert 2016] carries a distance-method
flag $f_D$ whose first value is defined in the machine-readable header as
\emph{``Hubble-Flow assuming $H_0 = 73$ km\,s$^{-1}$\,Mpc$^{-1}$ and
correcting for Virgo-centric infall''}. All counts below are on the standard
BTFR sample --- quality flag $Q \leq 2$ with a measured $V_{\rm flat}$,
which is $129$ of the $175$ galaxies in the table, not on the full $175$.

Not every non-flow method is equally free of $H_0$, and the distinction
matters enough to be made explicit (Table~\ref{tab:fD}). We separate three
tiers. \emph{Primary} distances (TRGB, Cepheids, SNe Ia; $33$ galaxies) are
measured to the individual galaxy by a stellar or standard-candle indicator:
these are redshift-independent in the NED-D sense. They are not free of the
distance-ladder \emph{zero point} --- TRGB and Cepheid calibrations share
anchors with SH0ES (NGC 4258, LMC eclipsing binaries, Gaia parallaxes) ---
but no value of $H_0$ enters the inference. The \emph{cluster} tier (Ursa
Major, $25$ galaxies) is ambiguous: SPARC assigns all $25$ members the
\emph{same} distance, $18.0$ Mpc, so it is a membership assignment rather
than an individual measurement, and the cluster distance itself rests on a
Tully--Fisher calibration that carries a distance-scale dependence. The
\emph{flow} tier ($71$ galaxies) is unambiguously $H_0$-dependent.

\begin{table}[h]
\centering
\begin{tabular}{llrr}
\toprule
Tier & Method & $N$ & fraction \\
\midrule
flow    & Hubble flow ($D = v/H_0$ at $H_0 = 73$) & 71 & $55.0\%$ \\
cluster & Ursa Major (all assigned $18.0$ Mpc)    & 25 & $19.4\%$ \\
primary & Tip of the red giant branch             & 28 & $21.7\%$ \\
primary & Cepheids                                &  3 & $2.3\%$ \\
primary & Supernovae Ia                           &  2 & $1.6\%$ \\
\midrule
\multicolumn{2}{l}{$f$ (loose: flow only)}        & \textbf{71} & $\mathbf{55.0\%}$ \\
\multicolumn{2}{l}{$f$ (strict: flow + cluster)}  & \textbf{96} & $\mathbf{74.4\%}$ \\
\multicolumn{2}{l}{genuinely redshift-independent} & \textbf{33} & $25.6\%$ \\
\bottomrule
\end{tabular}
\caption{Distance provenance of the standard SPARC BTFR sample ($129$
galaxies). Depending on how the Ursa Major cluster tier is counted, between
$55\%$ and $74\%$ of the sample carries distances that assume a value of
$H_0$ --- and that assumed value ($73$) sits at the SH0ES end of the
tension. Only $33$ galaxies have individually measured redshift-independent
distances.}
\label{tab:fD}
\end{table}

We carry both accountings forward, writing $f$ for the $H_0$-dependent
fraction: $f = 0.550$ (loose) or $f = 0.744$ (strict). Nothing below depends
on which is adopted.

For a galaxy whose distance comes from the Hubble flow, $D \propto 1/H_0$.
Baryonic masses are inferred photometrically, so $M \propto D^2$, while
$V_{\rm flat}$ is a velocity and carries no distance dependence. Since
$a_0 = V_{\rm flat}^4/(GM)$ in the BTFR normalisation,
\begin{equation}
a_0 \;\propto\; \frac{1}{M} \;\propto\; \frac{1}{D^2}
\;\propto\; \left(H_0^{\rm assumed}\right)^{2}
\qquad \text{on the } f_D = 1 \text{ sub-sample,}
\label{eq:coupling}
\end{equation}
The exponent in (\ref{eq:coupling}) is exact rather than estimated, and it
is worth seeing why. In the deep-MOND regime the interpolant gives
$a^2 = a_0 a_N$, so that $V^4/r^2 = a_0 GM/r^2$ and \emph{the radius cancels
identically}: the BTFR contains no length scale. The distance therefore
enters through the baryonic mass alone --- $V_{\rm flat}$ being a Doppler
velocity that carries no distance dependence --- and does so with the single
power $M \propto D^2$. Rescaling every distance in the sample by a common
factor $\lambda$ and recomputing returns $a_0 \propto \lambda^{-2}$ to five
decimal places, as it must.

A second structural observation sharpens what this does and does not
indict. In $M = V^4/(G a_0)$ the powers of time cancel exactly
($[V^4] = \rm m^4\,s^{-4}$, $[G a_0] = \rm m^4\,kg^{-1}\,s^{-4}$), so the
BTFR carries no net time dimension: $H_0$, an inverse time, cannot enter the
relation \emph{structurally}. It can enter only through the calibration of
the data --- the conversion of fluxes into masses, which requires a
distance. The circularity documented in this section is therefore a property
of the present SPARC distance mix, not of Eq.~(\ref{eq:law}), and that is
why the remedy is a better-selected sample rather than abandonment of the
relation.

Propagating (\ref{eq:coupling}) gives a sample-weighted sensitivity
$\mathrm{d}\ln a_0 / \mathrm{d}\ln H_0 = 2f$, that is $1.10$ (loose) or
$1.49$ (strict). Either way the coefficient exceeds unity: a $1\%$ error in
the assumed $H_0$ moves $a_0$ by more than $1\%$. A shift of the assumed
$H_0$ across the Planck--SH0ES range ($73 \to 67.4$, i.e.\ $-7.7\%$) moves
$a_0$ by $-8\%$ to $-11\%$: the same size as the tension the relation is
meant to probe.

The consequence is that Eq.~(\ref{eq:law}) and the calibration of $a_0$ are
not independent equations. Imposing both simultaneously,
\begin{equation}
H_0 \;=\; K\,g(H_0), \qquad
g(H) \equiv f\left(\frac{H}{73}\right)^{2} + (1-f), \qquad
K \equiv \frac{\pi\sqrt{3}\,a_0^{\rm SPARC}}{c} = 67.21,
\label{eq:fixedpoint}
\end{equation}
gives roots at $H_0 = 43.1$ and $100.9$ km\,s$^{-1}$\,Mpc$^{-1}$ for
$f = 0.550$, and at $21.6$ and $85.0$ for $f = 0.744$. None of the four is
$67.2$. The value quoted in \S3 is the solution of (\ref{eq:law}) at
\emph{fixed} $a_0$ --- that is, holding constant the very quantity that was
calibrated against $H_0 = 73$.

The conditioning of (\ref{eq:fixedpoint}) is more informative than its
roots. Differentiating, the response of the inferred $H_0$ to the published
$a_0$ is
\begin{equation}
\frac{\mathrm{d}\ln H_0}{\mathrm{d}\ln a_0}
\;=\; \frac{K}{H_0}\,\frac{g(H_0)}{1 - K\,g'(H_0)},
\qquad g'(H) = \frac{2fH}{73^2}.
\label{eq:conditioning}
\end{equation}
For $f = 0.550$ the denominator vanishes at $H_0 = 72.0$ --- \emph{inside}
the range under dispute. The response is $+13.7$ at $H_0 = 67.2$, $+32.5$ at
$70$, and changes sign to $-68.5$ by $73$: the loop is singular precisely
where the tension lives, and determines nothing. For $f = 0.744$ the
singularity moves out to $53.3$ and the response is a finite $-3.4$ to
$-2.5$ --- still large, and still \emph{negative}, meaning a higher
published $a_0$ implies a \emph{lower} self-consistent $H_0$, the opposite
of the naive reading of (\ref{eq:law}).

Under either accounting the loop fails to determine $H_0$, and that failure
--- not any particular root --- is the result of this section. We do not
present (\ref{eq:fixedpoint}) as a competing determination: it uses a linear
sample-weighted propagation, whereas the published $a_0$ comes from a global
fit to the radial acceleration relation rather than from a median over
Table~1, and the Virgo-centric infall correction is not a rigid rescaling.
The roots are an order-of-magnitude statement about the size of the
coupling, and that size is large enough to void the independence claim.

We also attempted the direct test: recomputing $a_0$ separately on the two
sub-samples. The $58$ galaxies with $H_0$-independent distances give a median
$a_0$ that is $+9.5\%$ higher than the $71$ Hubble-flow galaxies. However,
binning by $V_{\rm flat}$ shows the offset does not hold a consistent sign
($+43\%$, $+18\%$, $-11\%$, $+1\%$ from the lowest to the highest mass bin),
so it is degenerate with mass-dependent selection between the sub-samples
and is \emph{not} a clean measurement of the $H_0$ bias. We report it as a
bound on what Table~1 alone can resolve, not as a detection.

The clean route is available and we state it as the next step rather than
attempt it here: re-derive $a_0$ from the $58$ galaxies with
$H_0$-independent distances, using the full rotation curves and mass
modelling rather than the summary table. That determination would be
genuinely ladder-independent, and Eq.~(\ref{eq:law}) applied to it would be
the measurement this letter currently cannot make.

\section{The orthogonal companion: $\sigma_\Phi(z)$ as a probe of $E(z)$}

The BTFR flow derived from the interpolant $\mu(x) = x/\sqrt{1+x^2}$
[Figueroa Torres 2026b, \S4] produces the intrinsic-scatter evolution
\begin{equation}
\sigma_\Phi(z) \;=\; \sigma_\Phi(0)\,
\left(\frac{H(z)}{H_0}\right)^{\beta'(0)}
\;=\; \sigma_\Phi(0)\, E(z)^{\beta'(0)},
\label{eq:sigmaPhi}
\end{equation}
where $E(z) = H(z)/H_0 = \sqrt{\Omega_m(1+z)^3 + \Omega_\Lambda}$ and
$\beta'(0) = -1$ for the standard interpolant of the theorem. The
key observation is that $E(z)$ depends only on the density parameters
$\Omega_m, \Omega_\Lambda$ and \emph{not} on $H_0$ itself: the identical
$E(2) = 3.032$ is obtained for any $H_0 \in [65, 75]$ km\,s$^{-1}$\,Mpc$^{-1}$
holding $\Omega_m = 0.315$ fixed. Equation (\ref{eq:sigmaPhi}) therefore
measures the \emph{shape} of the expansion history, not its overall scale.

\subsection*{Symmetry protection: why (\ref{eq:sigmaPhi}) is well-defined}

The reconstruction (\ref{eq:sigmaPhi}) is not automatic. If the interpolating
function's Taylor expansion at the deep-MOND fixed point had a non-vanishing
quadratic coefficient $c_2$, the flow would become non-Lipschitz and
$\sigma_\Phi(z)$ would collapse to zero exactly at finite flow time
[Figueroa Torres 2026b, \S5], destroying the smooth $E(z)^{\beta'(0)}$
dependence on which the reconstruction relies. The Appendix A.11 of that
paper formalises this: $c_2 = 0$ is a knife-edge condition requiring
symmetry protection, and the isotropy axiom of the theorem
[Figueroa Torres 2026b, \S2] provides it exactly. The condition is moreover
not imposed by hand: among the four simple projection gates whose
solid-angle fractions generate the candidate interpolants, only the tilt
gate $g\sin\theta > a_0\cos\theta$ yields a $\mu$ with vanishing quadratic
coefficient [Figueroa Torres 2026b, \S3]. The reconstruction
(\ref{eq:sigmaPhi}) therefore rests on a symmetry-protected structural
identity with a geometric origin, not on a phenomenological choice. Any competing MOND-like framework
whose interpolant does not enforce $c_2 = 0$ will not yield a smooth
$\sigma_\Phi(z)$ curve, and its inability to do so will be visible in
first-generation ELT/JWST-IFU measurements at $z \gtrsim 1$.

\subsection*{Orthogonality with $a_0$}

The two observables carry orthogonal information about the expansion
history:
\begin{center}
\begin{tabular}{lcl}
\toprule
Observable & Anchor epoch & What it measures \\
\midrule
$a_0$ (SPARC) & $z \simeq 0$ & $H_0$ (scale) via Eq.~(\ref{eq:law}) \\
$\sigma_\Phi(z)$ (ELT/JWST-IFU) & $z \simeq 1$--$2$ & $E(z) = H(z)/H_0$ (shape) via Eq.~(\ref{eq:sigmaPhi}) \\
\bottomrule
\end{tabular}
\end{center}
Combined, they reconstruct $H(z) = H_0 \cdot E(z)$ across the redshift range
accessible to resolved kinematics, without invoking the CMB or the distance
ladder at any step. The reconstruction is closed even without external
priors on $\Omega_m$: two measurements of $\sigma_\Phi$ at distinct
redshifts fix both the shape parameter and, through
$E(z_1)/E(z_2)$, provide a consistency check against the assumed background
cosmology.

\section{Combined reconstruction and registered predictions}

Table~\ref{tab:pred} lists the predictions of $\sigma_\Phi(z)$ for the three
surviving interpolants of [Figueroa Torres 2026b, Table 1], under two
cosmological scenarios: the Planck-consistent $\Lambda$CDM background
($\Omega_m = 0.315, H_0 = 67.4$) and a SH0ES-consistent variant
($\Omega_m = 0.30, H_0 = 73.5$). Because $E(z)$ is invariant under the
choice of $H_0$ at fixed $\Omega_m$, the two scenarios differ only through
the modest $\sim 0.015$ shift in $\Omega_m$ typically accompanying the
higher-$H_0$ solutions. We anchor $\sigma_\Phi(0) = 0.07$ dex following the
observed SPARC intrinsic scatter [Lelli et al.\ 2019].

\begin{table}[h]
\centering
\begin{tabular}{l c c c c c}
\toprule
Interpolant & $\beta'(0)$ &
$\sigma_\Phi(z{=}0.5)$ &
$\sigma_\Phi(z{=}1.0)$ &
$\sigma_\Phi(z{=}2.0)$ &
$\sigma_\Phi(z{=}4.0)$ \\
\midrule
\multicolumn{6}{l}{\emph{Planck-consistent} ($\Omega_m = 0.315$, $H_0 = 67.4$):} \\
Standard $\mu = x/\sqrt{1+x^2}$
  & $-1$    & $0.053$ & $0.039$ & $0.023$ & $0.011$ \\
$\hat\mu$ entropic
  & $-2$    & $0.040$ & $0.022$ & $0.008$ & $0.002$ \\
$\tanh(x)$
  & $-2/3$  & $0.058$ & $0.047$ & $0.033$ & $0.020$ \\
\midrule
\multicolumn{6}{l}{\emph{SH0ES-consistent} ($\Omega_m = 0.30$, $H_0 = 73.5$):} \\
Standard $\mu = x/\sqrt{1+x^2}$
  & $-1$    & $0.053$ & $0.040$ & $0.024$ & $0.011$ \\
$\hat\mu$ entropic
  & $-2$    & $0.041$ & $0.023$ & $0.008$ & $0.002$ \\
$\tanh(x)$
  & $-2/3$  & $0.059$ & $0.048$ & $0.034$ & $0.021$ \\
\bottomrule
\end{tabular}
\caption{Registered predictions of $\sigma_\Phi(z)$ [dex] under two
cosmological scenarios. The insensitivity to $H_0$ at fixed $\Omega_m$
(compare rows across the two blocks) is the orthogonality property of \S5.
The dominant discriminator is $\beta'(0)$: for a target error budget of
$\pm 0.010$ dex per redshift bin, the standard interpolant is separable from
$\hat\mu$ near $z \simeq 1.2$ and from $\tanh$ near $z \simeq 2.3$ [Figueroa
Torres 2026b, Fig.~1]. Values computed at fixed $\sigma_\Phi(0) = 0.07$ dex.
Registered before ELT/JWST-IFU deployment.}
\label{tab:pred}
\end{table}

The insensitivity of $\sigma_\Phi(z)$ predictions to the choice of $H_0$
under fixed $\Omega_m$ is the direct expression of the orthogonality
established in \S5: the shape measurement does not carry scale information,
and vice versa. What the shape measurement \emph{does} carry is
discriminating power between models that resolve the Hubble tension through
modifications of $E(z)$. Early-dark-energy models, for instance, modify
$E(z)$ at $z \gtrsim 1000$ but leave the late-time expansion essentially
untouched; models that modify late-time dark-energy dynamics predict
deviations at $z \lesssim 1$. A measurement of $\sigma_\Phi(z=1.5)$ at
$\pm 0.005$ dex precision --- within reach of ELT/HARMONI resolved
kinematics with SED-fit baryonic masses --- would separate these classes at
the $\gtrsim 3\sigma$ level, independently of any determination of the CMB
acoustic scale.

\section{Epistemological disclosure}

We are explicit about what has been fit, what has been derived, and what
remains open.

\emph{Fit to data:} none of the numbers in this letter. The prefactor
$\pi\sqrt{3}$ in Eq.~(\ref{eq:law}) is a pure geometric constant; the
factor $4\pi$ in Eq.~(\ref{eq:o6}) is the Gauss-theorem projection weight,
structural to the AQUAL embedding and independent of the choice of
interpolating function. The value of $a_0$ entering
Eq.~(\ref{eq:H0central}) is the SPARC-canonical determination. It is
\emph{not}, however, independent of $H_0$: as quantified in \S4, $55\%$ of
the sample carries Hubble-flow distances computed at $H_0 = 73$, coupling
$a_0$ to the assumed expansion rate with sensitivity
$\mathrm{d}\ln a_0/\mathrm{d}\ln H_0 \simeq 1.1$. An earlier version of this
letter asserted that independence; the assertion was wrong and is retracted
here.

\emph{Branch ambiguity:} Eq.~(\ref{eq:law}) presupposes that the
acceleration scale tracks the instantaneous expansion rate,
$a_0 \propto H(z)$. This is one of three readings of the framework's scale
relation. The covariant reading $a_0 \propto \sqrt{R(z)}$ is excluded
independently of the present data, since $R = 8\pi G(\rho - 3p)$ vanishes
exactly in the radiation era while the expansion $\theta = 3H$ does not
[Figueroa Torres 2026, \S7.6]. The third reading, $a_0 \propto \sqrt{\Lambda}$
(constant in $z$), is not excluded, and it would make $\Lambda$-EG
indistinguishable from MOND in redshift.

Each branch predicts a different $a_0$ at fixed $H_0$, so each inverts to a
different $H_0$ at fixed $a_0$. Evaluated at $H_0 = 67.4$ the three
predictions are $a_0^{\rm A} = 1.203$, $a_0^{\rm B} = 1.051$ and
$a_0^{\rm C} = 0.995$, in units of $10^{-10}\,\rm m\,s^{-2}$ [Figueroa
Torres 2026, \S7.6]. Since each branch has $a_0 \propto H_0$ at fixed
$\Omega_\Lambda$, inverting on the observed $a_0 = 1.20$ gives
\begin{equation}
H_0^{\rm branch} \;=\; 67.4 \times \frac{1.20}{a_0^{\rm branch}/10^{-10}}
\;=\; 67.2\ (\text{A}), \quad 77.0\ (\text{B}), \quad 81.3\ (\text{C}),
\label{eq:branches}
\end{equation}
an ambiguity of $21\%$, larger than the $8.4\%$ tension being probed. The
discrepancy of each branch is properly quoted in $a_0$, where the $20\%$
systematic lives: branch C misses the central value by
$(1.20-0.995)/0.24 = 0.85\sigma_{\rm sys}$ and branch B by
$0.62\sigma_{\rm sys}$. (Quoting branch C's offset against the
\emph{branch-A} band of $\pm 13.4$ would give a different and incorrect
$1.05\sigma$: that band is $\pi\sqrt{3}\,\sigma_{a_0}/c$, the propagation of
$\sigma_{a_0}$ through branch A's constant, and does not apply to a
displacement generated by a different branch.) We note that
$0.85\sigma_{\rm sys}$ is weak grounds for exclusion. The reader should
treat $67.2$ as conditional on the $a_0 \propto H(z)$ branch.

We stress that (\ref{eq:branches}) is a different calculation from the
self-consistency roots of \S4 and the two should not be conflated: the
former asks which $H_0$ reproduces the observed $a_0$ under each scale
relation, the latter asks which $H_0$ is consistent with a dataset
calibrated at $H_0 = 73$. Setting $f = 1$ in (\ref{eq:fixedpoint}), for
instance, returns $73^2/K = 79.3$, which is neither $81.3$ nor a branch
statement at all.

\emph{Structural inputs:} the framework paper adopts the AQUAL embedding
as its scalar sector; the derivation of AQUAL from a more fundamental
principle within the covariant action of $\Lambda$-EG remains as Open
Problem 6 [Figueroa Torres 2026, \S13]. This does not affect the present
letter, since Eq.~(\ref{eq:o6}) holds within AQUAL as a mathematical
consequence, but it means the ultimate justification of Eq.~(\ref{eq:law})
inherits whatever justification AQUAL receives from future work.

\emph{Observational inputs:} the SPARC value of $a_0$ carries a $\sim 20\%$
systematic uncertainty at present. This uncertainty, propagated linearly
through Eq.~(\ref{eq:law}), produces the wide $\pm 13.4$ km\,s$^{-1}$\,Mpc$^{-1}$
band on $H_0$. Statistical arbitration of the Hubble tension will therefore
require the SPARC systematic to be reduced to $\lesssim 5\%$; the pathway is
through improved stellar-mass modelling (Gaia parallax priors), tighter
inclination priors from resolved kinematics, and larger samples (extended
SPARC, LSST-era radial-acceleration relation catalogues). Reducing the
systematic is necessary but not sufficient: the distance provenance of \S4
must be addressed separately, since it is a bias rather than a variance and
does not shrink with sample size.

\emph{Ex-ante predictions:} the numerical entries of Table~\ref{tab:pred},
registered on the deposition date of this letter, will be compared against
ELT/HARMONI IFU measurements of resolved BTFR kinematics at
$z \simeq 1$--$3$ when they become available. Should the observed
$\sigma_\Phi(z)$ deviate from the standard-interpolant prediction under
either cosmological scenario, one or more of the following will be
falsified: the isotropy axiom of the theorem, the AQUAL embedding, the
$\Lambda$CDM background, or the assumption that $a_0$ is redshift-invariant
in the observer frame.

\emph{What we do not claim:} we do not claim to have resolved the Hubble
tension, and we do not claim a third independent determination of $H_0$.
Two things stand in the way, and we have quantified both: the width of the
SPARC systematic on $a_0$ precludes statistical arbitration at the level
demanded by the $5\sigma$ CMB--distance-ladder discrepancy, and the
distance provenance of the sample (\S4) prevents the present-epoch channel
from being ladder-independent at all; the second is a bias, and \S4 shows
the resulting loop is ill-conditioned in exactly the range under dispute.
We claim, more modestly: that
$\Lambda$-EG fixes a rigid algebraic relation between $a_0$ and $H_0$ with
no free parameters; that evaluating it on current data returns a value
$0.3\%$ from Planck, which is suggestive but cannot be counted as evidence
until the circularity is removed; that the route to removing it is concrete
and identified (the $33$ galaxies with individually measured
redshift-independent distances); and that the
orthogonal companion (\S5), being a ratio and therefore immune to the
distance problem, opens a program to reconstruct the \emph{shape} of $H(z)$
from galactic dynamics alone.

\section{Discussion and outlook}

The $a_0$--$H_0$ coincidence has been noticed in the MOND literature since
at least Milgrom (1999), who observed that $a_0 \sim cH_0$ in
order-of-magnitude terms and interpreted it through the de Sitter Unruh
temperature. Milgrom (2009) sharpened the connection considerably, deriving
the deep-MOND limit from space-time scale invariance and conjecturing that
local gravitational physics would take exactly the deep-MOND form in an
exact de Sitter universe, on the grounds that the isometry group of $dS_4$
and the symmetry group of the deep-MOND field equation are both $SO(4,1)$.
The present letter is agnostic on that mechanism: Eq.~(\ref{eq:law}) is used
here only as an algebraic relation between two measured scales.
Klinkhamer \& Kopp (2011) derived the exact relation
$a_0 = 4\pi (c/\hbar)\, k_B T_{\min}$ from entropic gravity with a minimum
temperature. The present letter differs from those in two respects.
First, the prefactor $\pi\sqrt{3}$ in Eq.~(\ref{eq:law}) is derived, not
observed: it follows from the specific de Sitter spectral geometry of the
framework paper [Figueroa Torres 2026] combined with the AQUAL embedding.
Second, and more importantly, the coincidence is complemented here by the
orthogonal $\sigma_\Phi(z)$ observable, which promotes a numerical
coincidence into a two-channel reconstruction program.

The path forward is observational. Four developments would sharpen this
letter's contribution to the tension debate, of which the first is
prerequisite to the rest:
(0) re-derivation of $a_0$ restricted to galaxies with redshift-independent
distances --- $33$ have individually measured ones in SPARC, $58$ if the
Ursa Major cluster tier is admitted (\S4) --- which is what would convert
Eq.~(\ref{eq:law}) from a consistency relation into a measurement;
(i) reduction of the SPARC systematic on $a_0$ to $\lesssim 5\%$ through
Gaia-anchored mass modelling on extended galaxy samples;
(ii) first-epoch measurements of $\sigma_\Phi(z)$ at $z \simeq 1$--$2$ with
ELT/HARMONI resolved kinematics, tested against the entries of
Table~\ref{tab:pred};
(iii) independent theoretical work on the AQUAL sector of the framework,
either sharpening or falsifying the derivation of (\ref{eq:o6}) from more
fundamental principles.

We stress the falsifiability of the path. Eq.~(\ref{eq:law}) is a rigid
prediction: any determination of $a_0$ on $H_0$-independent distances that
lies more than $\sim 3\sigma$ below (\ref{eq:a0sparc}) forces $\Lambda$-EG
into tension with Planck; any such determination more than $\sim 3\sigma$
above forces it into tension with SH0ES. The qualifier is essential --- on
the present distance mix the test is partly self-referential, and tightening
the error bars alone would not make it decisive. Neither outcome is
currently accessible with SPARC data quality, but both will be within a
decade. Eq.~(\ref{eq:sigmaPhi}) carries no such caveat, being a ratio in
which $H_0$ cancels identically: any observed $\sigma_\Phi(z=1.5) \notin [0.012, 0.040]$
dex (the union of the surviving-interpolant predictions of
Table~\ref{tab:pred}, computed at $\sigma_\Phi(0)=0.07$ dex from SPARC:
$\hat\mu$-entropic $0.0125$, standard $0.0296$, tanh $0.0394$)
falsifies the combined framework.

\section*{Data availability}

All numerical evaluations are reproducible from the archived scripts
{\tt H0\_from\_a0.py} (Eq.~\ref{eq:law} and Table~\ref{tab:H0}),
{\tt sigmaPhi\_predictions.py} (Table~\ref{tab:pred}) and
{\tt a0\_H0\_selfconsistency.py} (\S4, Table~\ref{tab:fD} and
Eq.~\ref{eq:fixedpoint}). All three are deterministic and require only the
Python standard library; the third additionally reads the machine-readable
SPARC Table~1 ({\tt SPARC\_Lelli2016c.mrt}), archived alongside them. The
symbolic verification of the fixed-point structure underlying
Eq.~(\ref{eq:sigmaPhi}) is in {\tt sensitivity.py} (Appendix A.11 of the
companion paper, reproduced here), which requires SymPy. All symbolic
identities used in \S2 and \S5 have been verified in the archived scripts of
the companion papers [Figueroa Torres 2026, 2026b].

\section*{Funding and competing interests}

This work received no external funding. The author declares no competing
interests.

\section*{References}

\begin{itemize}\itemsep2pt
\item Bekenstein, J., Milgrom, M.\ 1984, ApJ 286, 7
\item Figueroa Torres, J.\,P.\ 2026, \emph{Lambda EG: Emergent Gravitational
  Coherence and Galactic Dynamics Without Particle Dark Matter}, Zenodo
  preprint, doi:10.5281/zenodo.19784004
\item Figueroa Torres, J.\,P.\ 2026b, \emph{A geometric characterization of
  the standard MOND interpolating function: $\mu(x) = x/\sqrt{1+x^2}$ as
  the isotropic distribution of $\cot\theta$}, Zenodo preprint (ORCID
  0009-0005-5297-8777)
\item Klinkhamer, F.\,R., Kopp, M.\ 2011, Mod.\ Phys.\ Lett.\ A 26, 2783
  [arXiv:1104.2022]
\item Lelli, F., McGaugh, S., Schombert, J.\ 2016, AJ 152, 157
\item Lelli, F., McGaugh, S., Schombert, J., Pawlowski, M.\ 2019, ApJ 872, 25
\item McGaugh, S., Lelli, F., Schombert, J.\ 2016, Phys.\ Rev.\ Lett.\ 117,
  201101
\item Milgrom, M.\ 1999, Phys.\ Lett.\ A 253, 273 [arXiv:astro-ph/9805346]
\item Milgrom, M.\ 2009, ApJ 698, 1630, \emph{The MOND limit from space-time
  scale invariance} [arXiv:0810.4065]
\item Planck Collaboration 2020, A\&A 641, A6 [arXiv:1807.06209]
\item Riess, A.\,G.\ et al.\ 2024, ApJ 977, 120 [JWST-anchored SH0ES update]
\end{itemize}

\end{document}
